% page 236 du fac-simile (image p-235 du scan)
% bloc 165.1 x 242.0 mm ; marge gauche 8.0 mm
\pgc{-12.700mm}{0.000mm}{
\l{\cel{6}228}
\l{}
\l{}
\l{\sou{ikonometro}. (\sou{fotogr.})\cel{3}Vizilo per qua onu povas determinar la}
\l{\cel{3}disto fokala di la objektalo fotografala, apta donar (ye}
\l{\cel{3}dimensioni donita) la imajo de fotografoto o de fotograf-}
\l{\cel{3}otajo. - DEFIS.}
\l{}
\l{\sou{ikosaedro}. (\sou{geom.})\cel{2}Solido 20-edra. - DEFIS.}
\l{}
\l{\sou{iktero}.\cel{2}(\sou{patol}.) Morbo qua flavigas la pelo, qua devenas de}
\l{\cel{3}perturbeso di la sekreco e di la fluo di la bilo. -DEFIS.}
\l{}
\l{\sou{iktiokolo}.\cel{2}Substanco diafana, solida, preparita ek la nato-}
\l{\cel{3}veziko di la esturjoni, uzata por facar kapsuli gelatinatra,}
\l{\cel{3}o la gumo-tafto. - DEFI.}
\l{}
\l{\sou{iktiologio}.\cel{2}La parto di la natur-historio, qua traktas la fishi.}
\l{}
\l{\sou{iktiosauro}. Reptero fosila, kun fisho-vertebri. L.\cel{1}\sou{ichtyosaurus}.}
\l{\cel{3}- DEFIRS.}
\l{}
\l{\sou{iktuso}. (\sou{patol.})\cel{2}Afekto subita, qua frapas quale stroko.}
\l{\cel{3}(kom ex.: iktuso apoplexiala, epilepsiala).}
\l{}
\l{\sou{\textquotedbl{}iktus\textquotedbl{}}.\cel{2}(\sou{metriko}).\cel{2}Frapo qua perceptigas la tempo forta di}
\l{\cel{3}verso-pedo.}
\l{}
\l{\sou{-il-}.\cel{2}Sufixo qua signifikas \textquotedbl{}instrumento por...\textquotedbl{} (ne \textquotedbl{}machino}
\l{\cel{3}por...\textquotedbl{}).}
\l{}
\l{il(u).-\cel{2}Pronomo personala maskulala, persono triesma (uzata}
\l{\cel{3}kande onu volas emfazar la sexuo).}
\l{}
\l{\sou{ileono}. (\sou{anat.})\cel{2}La lasta e maxim longa parto di la intestino}
\l{\cel{3}tenua. - EFIS.}
\l{}
\l{\sou{ilexo}.\cel{2}(\sou{bot.)}\cel{2}Arboro di la familio \textquotedbl{}ilicinei\textquotedbl{}, sempre verda,}
\l{\cel{3}di qua la folii esas briloza e dornoza. - L.\cel{1}\sou{ilex aquifolium}.}
\l{}
\l{\sou{iliako}.\cel{2}(\sou{anat.}) La maxim granda ek la tri osti flankala. EFISL.}
\l{}
\l{\sou{iliterata}. Qua ne savas lektar.}
\l{}
\l{\sou{iluminar}. (\sou{trans.)}\cel{3}Ornar (manuskripto, libro imprimita) ye}
\l{\cel{3}miniaturi, graburi, e c. - DEFIRS.}
\l{}
\l{\sou{iluzionar}.\cel{2}(\sou{netrans.})\cel{2}Subisar ta eroro sensala o spiritala,}
\l{\cel{3}qua igas aceptar lo semblanta kom lo reala. - DEFIRS.}
\l{}
\l{\sou{-im-}.\cel{2}Sufixo qua signifikas nombro fracionala.}
\l{}
\l{\sou{imaginar}. (\sou{trans.})\cel{2}Reprezentar ulo en sua spirito. - EFIS.}
\l{}
\l{imaginara. (\sou{matem.})\cel{2}Simbolo algebrala, ek radikalo duesma-}
\l{\cel{3}grada pri nombro negativa. - DEFIRS.}
}
